
\documentclass{article}
\usepackage{colortbl}
\usepackage{makecell}
\usepackage{multirow}
\usepackage{supertabular}

\begin{document}

\newcounter{utterance}

\centering \large Interaction Transcript for game `eqbench', experiment `eq\_bench', episode 146 with reflexion{-}llama{-}3.1{-}8b{-}instruct.
\vspace{24pt}

{ \footnotesize  \setcounter{utterance}{1}
\setlength{\tabcolsep}{0pt}
\begin{supertabular}{c@{$\;$}|p{.15\linewidth}@{}p{.15\linewidth}p{.15\linewidth}p{.15\linewidth}p{.15\linewidth}p{.15\linewidth}}
    \# & \multicolumn{2}{c}{Player} && \multicolumn{2}{c}{Game Master} \\
    \hline

    \theutterance \stepcounter{utterance}  
    & & & \multicolumn{4}{p{0.6\linewidth}}{
        \cellcolor[rgb]{0.9,0.9,0.9}{
            \makecell[{{p{\linewidth}}}]{
                \texttt{\tiny{[P1$\langle$GM]}}
                \texttt{You have attempted this same task before and did not succeed. These are the plans you wrote after those attempts.} \\
\texttt{Trial \#1: New plan: } \\
\\ 
\texttt{Upon re{-}examining the dialogue, I realize that I focused too much on the surface{-}level insults and didn't consider the underlying emotions and power dynamics at play. I should have paid closer attention to the fact that Jackson is accusing Alejandro of betraying him by being with another man's woman, and Alejandro is countering by implying that Jackson is a hypocrite who couldn't keep his own partner.} \\
\\ 
\texttt{To improve my analysis, I will focus on identifying the key emotional triggers and power struggles in the dialogue. I will also consider the context and history between the two characters, as well as the tone and language used by each of them.} \\
\\ 
\texttt{In particular, I will pay closer attention to the following:} \\
\\ 
\texttt{* The fact that Jackson is accusing Alejandro of betraying him, and how this accusation is likely to make Alejandro feel defensive and possibly guilty.} \\
\texttt{* The fact that Alejandro is countering by implying that Jackson is a hypocrite, and how this counter{-}attack is likely to make Jackson feel embarrassed and possibly betrayed.} \\
\texttt{* The power dynamics at play in the dialogue, with Jackson trying to assert his moral superiority and Alejandro pushing back against this assertion.} \\
\\ 
\texttt{By considering these factors, I hope to provide a more nuanced and accurate analysis of the emotional responses of the characters in the dialogue.} \\
\\ 
\\ 
\\ 
\texttt{Your task is to predict the likely emotional responses of a character in this dialogue:} \\
\\ 
\texttt{Jackson: Sanchez. Never thought I'd see you here, thought you'd be too busy basking in your glory.} \\
\texttt{Alejandro: Riley, always the charmer. I could say the same about you. How's the knee?} \\
\texttt{Jackson: Better than your conscience, I reckon.} \\
\texttt{Alejandro: Ah, the ever{-}righteous Riley. You think you're the only one who's had to fight for everything?} \\
\texttt{Jackson: I don't recall stepping on others to get to the top.} \\
\texttt{Alejandro: And I don't recall asking for your opinion.} \\
\texttt{Jackson: You never did. Too busy making a spectacle of yourself.} \\
\texttt{Alejandro: And you're too busy playing the victim. } \\
\texttt{Jackson: I ain't the one parading around with another man's woman.} \\
\texttt{Alejandro: And I ain't the one who couldn't keep her.} \\
\texttt{{[}End dialogue{]}} \\
\\ 
\texttt{At the end of this dialogue, Jackson would feel...} \\
\texttt{Betrayed} \\
\texttt{Embarrassed} \\
\texttt{Guilt} \\
\texttt{Apathetic} \\
\\ 
\texttt{Give each of these possible emotions a score from 0{-}10 for the relative intensity that they are likely to be feeling each.} \\
\\ 
\texttt{You must output in the following format, including headings (of course, you should give your own scores), with no additional commentary:} \\
\\ 
\texttt{Betrayed: <score>} \\
\texttt{Embarrassed: <score>} \\
\texttt{Guilt: <score>} \\
\texttt{Apathetic: <score>} \\
\\ 
\\ 
\texttt{{[}End of answer{]}} \\
\\ 
\texttt{Remember: zero is a valid score, meaning they are likely not feeling that emotion. You must score at least one emotion > 0.} \\
\\ 
\texttt{Your answer:} \\
            }
        }
    }
     \\ \\

    \theutterance \stepcounter{utterance}  
    & \multicolumn{4}{p{0.6\linewidth}}{
        \cellcolor[rgb]{0.561,0.737,0.561}{
            \makecell[{{p{\linewidth}}}]{
                \texttt{\tiny{[P1$\rangle$GM]}}
                \texttt{Betrayed: 8} \\
\texttt{Embarrassed: 9} \\
\texttt{Guilt: 0} \\
\texttt{Apathetic: 0} \\
            }
        }
    }
    & & \\ \\

    \theutterance \stepcounter{utterance}  
    & & & \multicolumn{2}{p{0.3\linewidth}}{
        \cellcolor[rgb]{0.9,0.9,0.9}{
            \makecell[{{p{\linewidth}}}]{
                \texttt{\tiny{[GM$|$GM]}}
                \texttt{Betrayed: 8} \\
\texttt{Embarrassed: 9} \\
\texttt{Guilt: 0} \\
\texttt{Apathetic: 0} \\
            }
        }
    }
    & & \\ \\

    \theutterance \stepcounter{utterance}  
    & & & \multicolumn{2}{p{0.3\linewidth}}{
        \cellcolor[rgb]{0.9,0.9,0.9}{
            \makecell[{{p{\linewidth}}}]{
                \texttt{\tiny{[GM$|$GM]}}
                \texttt{Betrayed: 8} \\
\texttt{Embarrassed: 6} \\
\texttt{Guilt: 0} \\
\texttt{Apathetic: 0} \\
            }
        }
    }
    & & \\ \\

    \theutterance \stepcounter{utterance}  
    & & & \multicolumn{2}{p{0.3\linewidth}}{
        \cellcolor[rgb]{0.9,0.9,0.9}{
            \makecell[{{p{\linewidth}}}]{
                \texttt{\tiny{[GM$|$GM]}}
                \texttt{game\_result = LOSE} \\
            }
        }
    }
    & & \\ \\

    \theutterance \stepcounter{utterance}  
    & & \multicolumn{4}{p{0.6\linewidth}}{
        \cellcolor[rgb]{0.95,0.95,0.95}{
            \makecell[{{p{\linewidth}}}]{
                \texttt{\tiny{[P1$\leftrightarrow$P]}}
                \texttt{Post{-}episode Reflexion (not sent to GM):} \\
\texttt{New plan: } \\
\\ 
\texttt{Upon re{-}examining the dialogue, I realize that I focused too much on the surface{-}level insults and didn't consider the underlying emotions and power dynamics at play. I should have paid closer attention to the fact that Jackson is accusing Alejandro of betraying him by being with another man's woman, and Alejandro is countering by implying that Jackson is a hypocrite who couldn't keep his own partner.} \\
\\ 
\texttt{To improve my analysis, I will focus on identifying the key emotional triggers and power struggles in the dialogue. I will also consider the context and history between the two characters, as well as the tone and language used by each of them.} \\
\\ 
\texttt{In particular, I will pay closer attention to the following:} \\
\\ 
\texttt{* The fact that Jackson is accusing Alejandro of betraying him, and how this accusation is likely to make Alejandro feel defensive and possibly guilty.} \\
\texttt{* The fact that Alejandro is countering by implying that Jackson is a hypocrite, and how this counter{-}attack is likely to make Jackson feel embarrassed and possibly betrayed.} \\
\texttt{* The power dynamics at play in the dialogue, with Jackson trying to assert his moral superiority and Alejandro pushing back against this assertion.} \\
\\ 
\texttt{To further improve my analysis, I will also consider the following:} \\
\\ 
\texttt{* The fact that Jackson's accusation is not just about Alejandro's actions, but also about his perceived moral failing.} \\
\texttt{* The fact that Alejandro's counter{-}attack is not just about defending himself, but also about challenging Jackson's moral authority.} \\
\texttt{* The fact that the dialogue is not just about the characters' emotions, but also about their relationships and power dynamics.} \\
\\ 
\texttt{By considering these factors, I hope to provide a more nuanced and accurate analysis of the emotional responses of the characters in the dialogue.} \\
            }
        }
    }
    & \\ \\

\end{supertabular}
}

\end{document}
